\section{The uncapped equal-magnitude nuisance class}
Suppose $p=1$ and $z_i\in\{-c,0,c\}$ for some $c>0$. The original uncapped relative-margin rule converges from every positive initial design. This special statement does not extend to the unequal-magnitude counterexample above.

The finite-refresh branch was proved earlier. With infinitely many refreshes, the epoch argument gives $q_k\to\bar v>0$, and the lower bound on fresh steps gives $\sum_k t_k=\infty$, with $0<t_k\le1/2$.
For real $b,u$, $q>0$ and $0<t\le1/2$,
\begin{equation}\label{eq:pair-product}
 [1-t+t(b+u)^2/q][1-t+t(b-u)^2/q]\ge1+2t(b^2/q-1).
\end{equation}
Indeed, put $\beta=b^2/q$, $x=u^2/q$. The product is
$F(x)=(1-t+t\beta+tx)^2-4t^2\beta x$.
If its unconstrained minimizer $x_*=\beta-(1-t)/t$ is nonpositive, its minimum on $x\ge0$ is $(1+t(\beta-1))^2$, proving the inequality.
Otherwise its minimum is $4t\beta(1-t)$; subtracting the right side gives
$(1-2t)(2t\beta-1)\ge0$, since $\beta>(1-t)/t$.

Orient nonzero-nuisance rows by $\operatorname{sign}(z_i)$ and put
\[
 b_i=\operatorname{sign}(z_i)a_i,\quad
 D_*=\tfrac12(\max_{z_i\ne0}b_i-\min_{z_i\ne0}b_i),\quad
 A_0=\max_{z_i=0}a_i^2,
\]
using $A_0=0$ if that set is empty. Then $v_*=\max\{A_0,D_*^2\}$:
choose the coefficient minus the midpoint of the extreme $b_i$, divided by $c$, for an upper bound; a zero-nuisance singleton attains $A_0$ when present. If $D_*>0$, the distinct extreme pair with equal masses attains $D_*^2$; if $D_*=0$, its lower bound is trivial, and no pair (or self-pair) is required.
Only in the case $\bar v<D_*^2$ is a pair argument needed; then $D_*>0$. After negating one residual in that pair, its residuals are $D_*+u_k,D_*-u_k$ for every cache, even if fits oscillate.
If $\bar v<D_*^2$, eventually $D_*^2/q_k\ge1+\eta$ for some $\eta>0$.
Equation~\eqref{eq:pair-product} multiplies their positive weight product by at least $1+2\eta t_k$. Its logarithm would diverge, since
$\log(1+2\eta t_k)\ge2\eta t_k/(1+\eta)$, contradicting the product bound $1/4$.
A zero-nuisance row similarly excludes $\bar v<a_i^2$ through growth of its weight.
Hence $\bar v\ge\max\{A_0,D_*^2\}=v_*$, proving equality without a weight floor or convergence of cached fits.

\section{A sharp high-accuracy penalty of the refresh-count cap}
Consider the capped $K$ rule on $(1,-1),(1,1)$ with first current mass $x=(1-y)/2$ and cache mass $s_1=(1-a)/2$, where $0<y\le a<1$. The cached residuals are $(1+a,1-a)$, $A_s=1$, and substitution gives
\[
 q=1+a^2-2ay,\quad d_1=-d_2=d=\frac{a(1-y^2)}q,\quad g=4ad,
\]
\[
 r=y-a,\quad h=-2d,\quad\alpha=\frac{1+y}{1+a},\quad
 \gamma=\frac{2d}{1-y},\quad\zeta=\frac{2d}{1+y},
\]
\[
 \frac Gg=1-\frac{1+a}{1+y}
 \left(\frac{a-y}{a}+\frac{(a-y)^2}{2a(1+y)}\right),\quad
 K=\frac{4a^2(1+y)(1+a)(1+a-2y)^2}{q^2}.
\]
For $z=y/a$, call the first ratio $F_a(z)$. The gate satisfies
\[
 2(1+az)^2(F_a(z)-1/2)
 =(a+2a^2)z^2+(2+4a)z-(1+3a+a^2)=:P_a(z).
\]
On $[0,1]$, $P_a'>0$, $P_a(1/2)=-3a/4-a^2/2<0$ and $P_a(1)=(1+a)^2>0$.
Its unique zero $z_*(a)\in(1/2,1)$ is the exact retained-cache threshold, with $z_*(a)=1/2+3a/8+O(a^2)$.

The true directional derivative is $4yd$, so $G\le4yd$.
Also $K=(2d+\gamma(a-y))^2/\alpha\ge4d^2$. Therefore
\[
 t\le G/(4K)\le y/(4d),\qquad y/2\le y^+=y-2td<y.
\]
No finite step hits or crosses the optimum. Capped convergence gives $y\to0$.
There must be infinitely many refreshes: a fixed final positive $a$ would contradict $y/a\ge z_*(a)>1/2$.
Let $a_j$ be the cache bias after refresh $j$, including the initial $j=0$, and $n_j$ its number of accepted steps.
Uniformly for $z\in[0,1]$, $g/a^2\to4$ and $K/a^2\to4$ as $a\to0$.
On the accepted set $G/g\ge1/2$, the uncapped safe step therefore has a fixed positive lower bound in all sufficiently late epochs. The cap $1/(j+1)$ eventually determines every step in those epochs.
Their exact recurrence is
\[
 z^+=z-\frac2{j+1}
       \frac{1-a_j^2z^2}{1+a_j^2-2a_j^2z}.
\]
The last factor converges uniformly to one and the threshold tends to $1/2$.
Exit overshoots it by at most $O(1/j)$, giving
\[
 z_{\rm exit}\to1/2,\qquad n_j/(j+1)\to1/4,\qquad
 a_{j+1}/a_j=z_{\rm exit}\to1/2.
\]
At the start of epoch $J$, $N_J=\sum_{j<J}n_j$ satisfies
$N_J/J^2\to1/8$ and $\log a_J/J\to-\log2$.
Within that epoch $N=N_J+O(J)$ and $\log y=\log a_J+O(1)$, because $y/a_J$ stays bounded away from zero.
The efficiency gap is $e_N=y_N^2$, so
\[
 \frac{\log e_N}{\sqrt N}\longrightarrow-4\sqrt2\log2.
\]
Strictly decreasing gaps permit inversion:
\[
 N_\epsilon\sim\frac{\log^2(1/\epsilon)}{32\log^2 2},\qquad
 J_\epsilon\sim\frac{\log(1/\epsilon)}{2\log2}.
\]
For the uncapped rule every fresh state $y=a$ has $d=a$, $g=G=K=4a^2$, so $t=1/4$ and $y^+=a/2$.
The strict inequality $P_a(1/2)<0$ forces a refresh before the next step.
Thus $y_N=2^{-N}y_0$ exactly.
Fresh half-affine averages the swapped masses, while half-power gives equal unnormalized masses $2x(1-x)$; both are optimal in one update.
Repeating the $+1$ row gives a full-rank uniform three-row initialization with aggregate $x_0=1/3$.
This is a duplicate-row example, not a distinct-site uniform-start result.

\paragraph{Generic-operation translation.}
For a completed nonstationary prefix of $N\ge1$ accepted steps and $J$ actual refreshes, stop counting at the final accepted update.
There are $N+J$ gates, of which $J+1$ are fresh and $N-1$ stale.
The recorded generic implementation consequently has
\[
 F=J+1,\quad H=2N+J-1,\quad M=2N+3J+1,\quad T=2(N+J).
\]
Here $F$ counts each assembly, factorization and cache contrast solve; $H$ old-factor solve pairs; $M$ feature passes; $T$ scalar reduction passes.
Fresh gates pay one old solve and stale gates two, including rejected gates; each cache also pays two feature passes.
Let fixed nonnegative prices be $B_0,h,\mu,\nu$, with $B_0$ the sum of assembly/factor/contrast prices. Then $C_K=B_0F+hH+\mu M+\nu T$.
Writing $L_\epsilon=\log(1/\epsilon)$, for $h+\mu+\nu>0$,
\[
 C_{\rm cap}\sim\frac{h+\mu+\nu}{16\log^2 2}L_\epsilon^2,\qquad
 C_{\rm uncapped}\sim\frac{B_0+3h+5\mu+4\nu}{2\log2}L_\epsilon.
\]
Their cost ratio diverges, while the leading refresh counts agree.
These are operation prices, not measured hardware coefficients.
Special scalar formulas may remove solves for both implementations, but not the accepted-update count separation.
Additional priced nuisance inner products are considered separately below; no hardware speedup follows.

\section{First-hit cost regions including preconditioned correction}
This is a static calculation from saved runs, with no new trajectory, timing or LP execution. Every row uses that method's first recorded offline efficiency at least $0.999999$, not a later matched-budget iterate. The target is not an online decision input.
Extend the preceding price model to
\[
 C=B_0F+hH+\mu M+\nu T+\xi D,
\]
where $D$ counts explicit nuisance-vector inner products outside solves and feature matvecs. These length-$p$ dots differ from length-$m$ node reductions. This symbolic model still does not measure all physical FLOPs, memory traffic or hardware overhead.
Old $K$ and $B_{\rm sur}$ gates have $D=H$. A correction computes two dots for its line minimizer and one more only when its corrected residual is nonzero. If $N_c$ corrections and $N_y$ nonzero corrected residuals occur, $H=N_c+N_y$ and $D=2N_c+N_y=H+N_c$.
The saved primitive counts record no zero-initial-residual stale skips on the three runs, so the two attained correction first hits have $N_c=N-1$. Fresh half-power/half-affine have $D=0$. Dots internal to old-factor triangular solves are already paid in $H$; those internal to a fresh contrast solve are paid in $B_0$.

\begin{center}\small
\begin{tabular}{llrrrrrr}
\toprule
Case & Method & Steps & $F$ & $H$ & $M$ & $T$ & $D$\\
\midrule
$(1,2)$ & Half-power &1&1&0&2&1&0\\
 & Half-affine &2&2&0&4&2&0\\
 & Uncapped $K$ &8&8&22&38&30&22\\
 & Capped $K$ &9&6&22&34&28&22\\
 & Capped $B_{\rm sur}$ &10&6&24&36&45&24\\
 & Correction &9&1&13&26&27&21\\
\midrule
Cubic & Half-power &203&203&0&406&203&0\\
 & Half-affine &206&206&0&412&206&0\\
 & Uncapped $K$ &210&106&524&736&630&524\\
 & Correction &879&12&1756&2658&2670&2634\\
\bottomrule
\end{tabular}
\end{center}
The old-fit capped cubic methods miss the target, and all six methods miss it on the difficult boundary case; those rows are excluded from first-hit comparisons. The $H$ column is old solves, not the total-solve column in the main experimental table.

Half-power is the unique Pareto-efficient method on $(1,2)$; correction itself dominates all three old-fit safeguards there. For the cubic case, half-power dominates half-affine, leaving half-power $P$, uncapped $U$ and correction $C$. Define
\[
 T_{UP}=\frac{524h+330\mu+427\nu+524\xi}{97},\qquad
 T_{CU}=\frac{1232h+1922\mu+2040\nu+2110\xi}{94}.
\]
Every coefficient of $T_{CU}$ exceeds its counterpart in $T_{UP}$. For strictly positive prices $T_{UP}<T_{CU}$, and the correction/half-power threshold is
\[
 T_{CP}=\frac{97T_{UP}+94T_{CU}}{191}.
\]
Subtracting the table's resource vectors proves that $P$ is cheapest for $B_0<T_{UP}$, $U$ for $T_{UP}<B_0<T_{CU}$, and $C$ for $B_0>T_{CU}$, with ties at the boundaries. For example, correction beats half-power precisely when
$191B_0>1756h+2252\mu+2467\nu+2634\xi$.
Setting $\xi=0$ recovers the four-price projection; it deliberately omits a separately priced resource rather than claiming exhaustive arithmetic accounting.

For generic $K$ prefixes the extension replaces $h$ by $h+\xi$ in the preceding cost asymptotics, since $D=H$. Thus its logarithmic-squared separation persists for fixed nonnegative prices with $h+\xi+\mu+\nu>0$. It does not follow if only cache construction is priced, or if prices change with requested accuracy.
All first-hit determinations use separately charged offline truth. No certified online stop or physical price calibration is added by this calculation.

\paragraph{A direct inexact-gradient comparison.}
On this cap-cost instance the late epoch lengths satisfy $n_j/(j+1)\to1/4$, and each update of that epoch has $t_k=1/(j+1)$. Therefore
\[
 \sum_{k\text{ in epoch }j}t_k^2
 =\frac{n_j}{(j+1)^2}\sim\frac1{4(j+1)}.
\]
The sum over all accepted steps diverges (eventually each epoch contributes at least $1/[8(j+1)]$). Thus identifying $\gamma_k=t_k$ does not meet the square-summability hypothesis of Bertsekas and Tsitsiklis~\cite[Proposition~1]{bertsekas2000}. That proposition additionally assumes a globally Lipschitz gradient on an unconstrained domain and uniform direction/error bounds; these are not supplied by the simplex slope gate alone. The comparison does not rule out other transformations or convergence theorems.
